Single-cell foundation models (scFMs) such as Geneformer~\cite{theodoris2023transfer} and
scGPT~\cite{cui2024scgpt} are used for cell-type annotation, perturbation-response prediction, and
gene-network inference. They are trained on millions of transcriptomic profiles and learn
contextual gene representations without explicit supervision on regulatory relationships. A central
question is whether these representations encode \emph{causal regulatory logic}---the directed
relationships between transcription factors (TFs) and their target genes---or only the
\emph{statistical co-expression} that correlates with, but does not constitute, regulation.

At the level of attention weights this question has been addressed by a companion
study~\cite{kendiukhov2025systematic}, which found that attention captures co-expression rather
than regulatory signal: trivial gene-level baselines outperform attention edges, pairwise scores
add zero incremental predictive value, and causal ablation of putatively regulatory heads produces
no behavioural effect. Attention weights, however, are one view among several. The residual
stream---the running sum of all layer outputs---may encode richer structure than attention patterns
alone reveal.

The superposition hypothesis~\cite{elhage2022superposition} offers a framework for this richer
structure. When a model must represent more concepts than it has dimensions, it encodes features as
nearly-orthogonal directions in activation space and activates only a sparse subset of them per
input. Sparse autoencoders
(SAEs)~\cite{olshausen1997sparse,sharkey2022taking,cunningham2023sparse,bricken2023monosemanticity}
resolve superposition by learning an overcomplete dictionary with a sparsity constraint,
decomposing dense activations into interpretable features. In large language models, SAE pipelines
have scaled to billions of features~\cite{gao2024scaling,templeton2024scaling}; they have not been
systematically applied to biological foundation models, where the ``concepts'' are genes, pathways,
and regulatory programs rather than linguistic entities.

Here we present a systematic SAE-based interpretability analysis of single-cell foundation models.
We train TopK SAEs~\cite{makhzani2013ksparse} on per-position residual-stream activations from all
18 layers of Geneformer V2-316M and all 12 layers of scGPT whole-human~\cite{cui2024scgpt},
producing atlases of 82,525 and 24,527 features respectively. The two models differ substantially
in architecture (Geneformer's rank-value tokenization vs.\ scGPT's dual gene-ID plus binned-value
embeddings, 18 vs.\ 12 layers, $d{=}1{,}152$ vs.\ $d{=}512$, bidirectional masked-token vs.\
specialized causal-attention generative training) and training data (30M vs.\ 33M cells), yet we
apply an identical analytical pipeline to both.

Two questions organize the analysis. The first is descriptive: what biological structure do these
dictionaries contain, how is it organized across layers, and how much of it is accessible to
standard linear decomposition of the same activations? The second is causal: do the features
encode which genes a transcription factor actually regulates? Answering the second question
quantitatively requires a reference point, because a low recovery rate is uninformative unless one
knows how much of the regulatory relationship is recoverable from the same measurements by any
method. We therefore benchmark the SAE against a perturbation-aware empirical ceiling---
differential expression of the CRISPRi knockdown itself---and against established
regulatory-inference baselines evaluated in an identical framework, and we replicate the
measurement in four cell lines.

To make these results accessible to the community, we release both atlases as interactive web
platforms---the Geneformer Feature Atlas
(\url{https://biodyn-ai.github.io/geneformer-atlas/}) and scGPT Feature Atlas
(\url{https://biodyn-ai.github.io/scgpt-atlas/})---enabling exploration of over 107,000 features
across 30 layers of two leading single-cell foundation models.
